% ======================================================================
%  Appendices of The Actualization Theory
%  TQM-MONO115 · Canonical Monograph (v2.0) · Zenodo-ready
%  Includes: Appendix A (Universality Program — future work),
%            Appendix B (References structure),
%            Appendix C (Validation targets summary).
%  Assembly only — no new physics, no new derivations.
% ======================================================================

\chapter{The Universality Program: Future Work}
\label{app:universality}

The universality research program of The Actualization Theory is treated as future work, outside
the core physics derivation of this monograph. The core monograph derives the physics of the
theory; the universality program investigates the operator basis across organized domains
and the organization structure of spectra. Its results are summarized here for the record,
with the understanding that they are a separate research direction, not a requirement for
the physics derivation.

\section{The Operator Basis}
\label{app:operators}

The four operators --- Crowding (degeneracy grouping), Compression (octave banding), Beat
(frequency ratio), and Locking (spectral gap) --- are the derived reading instruments of the
D96 spectrum. They are projections of the spectrum, not primitives. No fifth operator has
been found in any searched domain.

\section{The Lock Law}
\label{app:lock-law}

The lock identities of organized spectra carry a universal structure (the presence of stable
ratios) with domain-specific values (each domain's fingerprint). The lock structure
precedes organization, and organization emerges at a critical transition. These results are
established on synthetic deterministic evolving-law cohorts, with the synthetic basis
disclosed.

\section{Scope}
\label{app:universality-scope}

The universality program strengthens the theory but is not required for the physics
derivation of Chapters 1--12. The operator basis and the lock law are presented here as the
future-work appendix of the monograph, consistent with the physics-focused publication plan.

\chapter{References}
\label{app:references}

The references of this monograph follow the canonical research-program record. Each chapter
carries its own bibliography of the accepted TQM-QG phases that support its results. The
reference structure is:

\begin{itemize}
\item \textbf{Foundation chapters (1--2):} the primitive and reference phases
(QG268, QG278, QG279, QG286, QG289--292).
\item \textbf{Dynamics chapters (3--5):} the actualization, closure, and spectrum phases
(QG267--268, QG282, QG293--296).
\item \textbf{Spectrum chapters (6--8):} the D96, operator, and lock phases
(QG155--160, QG260--263, QG300--319).
\item \textbf{Physics chapters (9--12):} the quantum, gravity, standard-model, and
cosmology phases (QG216--244, QG149--180, QG227--240).
\item \textbf{Boundary and frontier chapters (13--14):} the boundary and prediction phases
(QG185, QG196, QG242, QG245, QG190--203, QG299).
\end{itemize}

The complete machine-readable record of the research program is the physics-coverage
register, which ships with the Zenodo deposit.

\chapter{Validation Targets Summary}
\label{app:validation}

The following derived predictions await experimental validation:

\begin{center}
\begin{tabular}{lll}
\hline
\textbf{Target} & \textbf{Derived value} & \textbf{Validation path} \\
\hline
Oblique $S,T,U$ & $S=0.0421$, $T=0.0842$, $U=0$ & EW global-fit re-analysis \\
Electron g-2 $a_e$ & $1.159655\times10^{-3}$ & lattice-limited measurement \\
Majorana $m_{\beta\beta}$ & $2.02\times10^{-3}$ eV & nEXO / LEGEND-1000 \\
P1 (106 GeV) & $106.39$ GeV, window 99--114 & HL-LHC \\
P2 ($0\nu\beta\beta$) & $m_{\beta\beta}=2.02$ meV & nEXO / LEGEND-1000 \\
P3 (sector ladder) & 9 resonances to 263.43 GeV & HL-LHC confirmation \\
\hline
\end{tabular}
\end{center}

Each target carries an explicit falsification condition, documented in the prediction
registry and the outcome dashboard of the research program.
